\section{Training}

In this section, the training phase of the model is described. Mean Squared Error (MSE) was chosen as loss function; other functions like MAE or Huber loss were tested but showed lower performance than MSE. $k-$fold Cross Validation ($k-$CV) was used to estimate the generalization performance and to mitigate overfitting. For each fold, the model was trained on $k-1$ subsets and validated on the remaining one, rotating through all folds. This approach ensures that each data point is used for both training and validation, providing a more robust evaluation of model performance. A $k$ value of 5 was used.

Each model was trained from scratch through the ADAM optimizer \cite{kingma2017adammethodstochasticoptimization} with the hyperparameters found in the previous section.

\subsection{ADAM}
Adam combines the strengths of two other popular optimizers: 
\begin{itemize}
    \item Stochastic Gradient Descent (SGD) with momentum:  updates the parameters by combining the current gradient with an exponential average of past gradients, so as to accelerate optimization in consistent directions and reduce oscillations.
    \item RMSProp: it dynamically adjusts the learning rate of each parameter by dividing the gradient by an average of its recent magnitudes, improving stability and convergence especially with noisy or differently scaled gradients.
\end{itemize}

This combination of techniques allows ADAM to be more efficient and robust on noisy datasets.

\subsection{Training results}

The results are shown in Table~\ref{tab:training}. A Pareto dominance of MobileNetV3 is evident, as it achieves, by far, the lowest loss values across all metrics compared to the other models. This indicates that MobileNetV3 outperforms, in terms of accuracy, both AlexNet and SqueezeNet. This superiority was to be expected since AlexNet is an obsolete model compared to MobileNetV3 and SqueezeNet focuses more on model reduction than accuracy. 

\begin{table}[H]
    \centering
    \begin{tabular}{|r|c|c|c|c|} \hline
        & \textbf{Train loss}   & \textbf{Valid. Loss}   & \textbf{Test Loss} & \textbf{$5-$CV Loss} \\\hline
        \textbf{AlexNet}     & 0.03394 & 0.03214 & 0.03476 & 0.03489 \\\hline
        \textbf{MobileNetV3} & \color{red}0.00707 & \color{red}0.01499 & \color{red}0.01450 & \color{red}0.01362  \\\hline
        \textbf{SqueezeNet}  & 0.03703 & 0.03447 & 0.03849 & 0.03682 \\\hline
    \end{tabular}
    \caption{Models performances\label{tab:training}}
\end{table}

\noindent
Regarding efficiency, in terms of inference time and model size (Table \ref{tab:training_efficency}), as expected, SqueezeNet is the most lightweight model, achieving the smallest size and fast inference, especially on GPU, maintaining similar accuracy to the much heavier AlexNet. \

While not as small as SqueezeNet, MobileNetV3 is significantly more compact and faster than AlexNet, and it achieves the best accuracy among all models.

Inference times were calculated by averaging the individual datapoint inference times over the testset.

\begin{table}[H]
    \centering
    \begin{tabular}{|r|c|c|c|} \hline
        & \textbf{tCPU(ms)}   & \textbf{tGPU(ms)}   & \textbf{Size(MB)} \\\hline
        \textbf{AlexNet}     & 5.3542 & 0.2395 & 142.29 \\\hline
        \textbf{MobileNetV3} & 3.3972 & 0.8319 & 5.85 \\\hline
        \textbf{SqueezeNet}  & 3.8984 & 0.3311 & 2.76 \\\hline
    \end{tabular}
    \caption{Models inference times with CPU and GPU and size\label{tab:training_efficency}}
\end{table}